\chapter{Problem Statement}

\section{Rationale}

High-frequency bus operation is a sequential control problem: a holding or stop-skipping decision changes the headways, queues, loads, and downstream dwell times encountered by later buses. Fixed timetables cannot adapt after these quantities depart from their planned values \cite{Ceder2007,Daganzo2009}. Reinforcement learning is therefore relevant because it learns an action policy from repeated state transitions rather than using a fixed dispatch rule.

The empirical case is CapMetro Rapid Route 801 in Austin, Texas. CapMetro's public APC dataset covers July--December 2021 and contains event timestamps, route and direction codes, stop IDs, boarding and alighting counts, onboard load, dwell time, revenue travel time and distance, event coordinates, and quality indicators \cite{TexasCapMetroAPC2021}. Under identical cleaning rules, the reproducible audit found 455{,}654 clean stop events for Route 801 and 376{,}801 for Route 803. Route 801 also had more recorded boardings and more positive-time/positive-distance segments, which justifies its selection as the data-richer primary case. The decision does not imply superior or inferior service quality.

The one-direction study subset uses direction code 6 and contains 229{,}421 clean stop events, 184 service-day codes, and 29 distinct stop IDs. The direction is not assigned a compass label because a checksum-verified 2021 GTFS snapshot has not yet been acquired. Historical stop names, route shapes, scheduled headways, and capacity values therefore remain gated rather than being borrowed from the current schedule.

NOAA LCDv2 observations make an empirical ordinary-weather exposure feasible \cite{NOAALCDv2}. After the documented local-standard-time conversion, the nearest-observation join matched all primary APC events at both Camp Mabry and Austin-Bergstrom within 90 minutes, including 11{,}804 Camp Mabry rain-exposed events. Join coverage alone does not establish causality. Ordinary-rain effects must be estimated within comparable segment, time-of-day, and day-type strata; severe weather outside the observed support remains an explicitly synthetic stress test. Vehicle breakdowns are also synthetic because the APC dataset contains no failure-event field.

The resulting problem is to determine whether a parameter-sharing MARL controller, calibrated to observed CapMetro demand, dwell, load, and segment travel-time behavior, remains competitive with fixed-rule controllers when controlled demand surges, weather stress, and breakdowns are introduced without misrepresenting those synthetic variables as measured CapMetro incidents.

\section{Research Gap}

Existing MARL bus-control studies report improvements in waiting time or headway regularity, but many experiments retain stationary demand, symmetric travel-time noise, or a failure-free fleet \cite{Wang2020Holding,Rodriguez2023Cooperative,Wangsun}. The literature review identified no MARL bus-scheduling study that combines demand surges, ordinary traffic variability, heavy-tailed weather stress, and discrete breakdowns in one public-data-calibrated corridor evaluation. This is a robustness-evaluation gap, not evidence that all four disturbances were observed together in Austin.

The CapMetro and NOAA sources narrow that gap in two ways. First, they replace assumed baseline demand, dwell, load, and segment travel-time profiles with reproducible public observations. Second, they permit observed ordinary-rain exposure to be separated from synthetic out-of-support stress. They do not supply passenger arrival times, capacity, breakdowns, or an authoritative historical schedule. The study therefore distinguishes measured calibration inputs, derived variables, simulated passenger states, and synthetic stressors in every claim.

\section{Research Objectives}

\subsection{Main Objective}

To evaluate a parameter-sharing Multi-Agent Reinforcement Learning scheduling model under baseline and disturbed operating conditions in a public-data-calibrated simulation of CapMetro Rapid Route 801 direction code 6.

\subsection{Specific Objectives}

\begin{enumerate}
    \item \textbf{Environment construction:} construct a reproducible Route 801 simulation from the cleaned CapMetro APC subset and temporally aligned NOAA weather observations.

    \textbf{Expected Output 1.1:} A documented acquisition and preprocessing pipeline with source URLs, queries, cleaning rules, row counts, checksums, field mappings, time-zone handling, and unresolved data gates.

    \textbf{Expected Output 1.2:} A corridor model validated against held-out empirical stop-event counts and segment travel times, with any schedule- or capacity-dependent value left as \%TODO-DATA until its source is verified.

    \textbf{Expected Output 1.3:} Configurable demand-surge, observed-weather, synthetic severe-weather, and synthetic breakdown components whose activation is explicit in each experiment.

    \item \textbf{Algorithm implementation:} develop a shared-parameter DDQN bus-control policy with per-agent local observations, per-agent rewards, a shared replay buffer, and decentralized execution.

    \textbf{Expected Output 2.1:} A fully specified observation, action, reward, transition, and training protocol, including reward-weight sensitivity analysis and reproducible random-seed control.

    \item \textbf{Performance evaluation:} compare MARL with No Control, Forward Headway, and Even Headway controllers under a declared baseline, combined-disturbance test, and single-disturbance ablations.

    \textbf{Expected Output 3.1:} Matched-seed estimates and uncertainty intervals for simulated passenger waiting time, bus travel time, and headway coefficient of variation under the baseline evaluation.

    \textbf{Expected Output 3.2:} Matched-seed robustness estimates under combined and isolated demand, weather, and breakdown stress, reported relative to the baseline without labeling synthetic intensities as observed meteorological categories.
\end{enumerate}

\section{Significance of the Study}

\textbf{Practical significance.} The work provides CapMetro and comparable transit researchers with a reproducible route-level evidence package rather than an undocumented synthetic baseline. Its route choice, data exclusions, NOAA join, and source gaps are independently auditable. If the controller improves simulated service, the result indicates where dynamic control warrants further study; it is not an on-road performance claim.

\textbf{Scientific significance.} The study separates three layers that are often conflated: empirical ordinary operation, empirically joined weather exposure, and synthetic out-of-support stress. This separation enables a clearer robustness claim. The evaluation also uses matched random seeds across controllers and reports degradation relative to an observed-data-calibrated baseline, allowing later algorithms to be compared under the same corridor and disturbance definitions.

\section{Scope and Limitations}

\textbf{Scope.} The empirical scope is CapMetro Route 801, direction code 6, during July--December 2021. The simulation uses 29 observed stop IDs and a single simulated operating day per episode. The active fleet size, service horizon, scheduled headway, capacity, and control-stop set are implementation parameters that must be derived from verified sources or retained as \%TODO-DATA. Each active bus is an agent that uses the same learned DDQN weights but produces its own action from a local observation.

The experimental activation matrix is fixed as follows. Baseline stochastic demand (D) and ordinary empirical segment travel-time variation (T) are present in training and every evaluation. Training domain-randomizes demand surge (S), weather stress (W), and breakdown (B). Stage A evaluates D+T only, with S/W/B disabled. Stage B evaluates D+T+S+W+B. Three ablations add S only, W only, or B only to the Stage A baseline. Thus W never replaces T, and an ablation labeled ``weather only'' means D+T+W rather than a deterministic corridor with weather as its sole source of randomness.

Ordinary observed rain is joined from Camp Mabry, with Austin-Bergstrom used for station sensitivity. Temperature, relative humidity, wind, and visibility may enter as control covariates, but only precipitation and rain/weather codes define the primary observed exposure. A lognormal coefficient-of-variation sweep may be used only for severe/extreme stress outside the empirical support and must be labeled synthetic. The breakdown process is likewise synthetic and must be reported separately from observed APC facts.

The MARL policy is compared with No Control, Forward Headway, and Even Headway using simulated passenger waiting time, bus travel time, and headway coefficient of variation over at least 30 matched-seed Monte Carlo runs per evaluation cell.

\textbf{Limitations and delimitations.} (a) APC boardings do not identify passenger arrival times, so passenger waiting time is generated by the simulation and cannot be called directly observed. (b) APC coordinates are stop-event points, not continuous trajectories; calibration therefore uses segment travel time and event-level location quality rather than claiming continuous speed trajectories. (c) The units of \texttt{rev\_distance} depend on external odometer metadata and remain unresolved; travel-time RMSE is primary until units are verified. (d) Capacity and 2021 schedule semantics require separate authoritative sources. (e) Direction code 6 remains unlabeled until a compatible historical GTFS snapshot is checksum-verified. (f) No breakdown is inferred from missing or delayed APC records. (g) Feeder-road traffic is not simulated explicitly because the agent observes route-level headway, load, and queue states; its effect enters only through empirical demand and segment travel-time distributions. (h) No on-road deployment or sim-to-real claim is made.
